% !TeX root = ../main.tex
% !TeX encoding = utf8

\chapter{Ajuste estructural y conclusión}
% Responsable: José Ángel Carretero Montes
% Secciones del índice:
%   5. El ajuste estructural a los factores de contingencia
%   6. Conclusión

\section{Ajuste estructural a los factores de contingencia}

% Identificar y describir los factores de contingencia que afectan a CaixaBank
% y analizar si existe ajuste estructural.

\subsection{Edad y tamaño}

% ¿Cómo influyen la antigüedad de la entidad y su gran tamaño en su estructura?
%   - ¿Se aprecia mayor formalización respecto a entidades más jóvenes o pequeñas?
%   - ¿Qué consecuencias tiene la escala de CaixaBank sobre sus procedimientos?

\subsection{Sistema técnico}

% ¿Qué herramientas tecnológicas condicionan la forma de trabajar?
%   - ¿Hasta qué punto regulan o limitan la actividad operativa?
%   - ¿Se han automatizado tareas anteriormente manuales? ¿Con qué impacto en plantilla?

\subsection{Entorno}

% ¿Cómo es el entorno competitivo de CaixaBank?
%   - Estabilidad/dinamismo: ¿predecible o volátil (fintech, regulación, tipos de interés)?
%   - Complejidad: ¿los productos requieren conocimiento especializado?
%   - Cambios recientes en la organización derivados de cambios del entorno.

\subsection{Poder}

% ¿En qué medida agentes externos determinan la estructura y procedimientos de CaixaBank?
%   - Supervisión del BCE y el Banco de España.
%   - Normativa europea: MiFID II, Basilea III, DORA, etc.
%   - ¿Limitan la capacidad de decisión interna de la organización?

\subsection{Análisis del ajuste estructural}

% ¿Existe coherencia entre los factores de contingencia identificados y la estructura
% organizativa descrita en los capítulos anteriores?
%   - Si hay ajuste: argumentar por qué la estructura responde adecuadamente a los factores.
%   - Si no hay ajuste: aportar una razón que explique la desviación.

\section{Conclusión}

% Identificar a qué configuración estructural se asemeja más CaixaBank:
%   - Estructura simple
%   - Burocracia maquinal
%   - Burocracia profesional
%   - Forma divisional
%   - Adhocracia

% Argumentar la decisión indicando:
%   a) Características de CaixaBank que son COHERENTES con la configuración identificada.
%   b) Características que NO son coherentes con dicha configuración.

% Valorar si se trata de una estructura híbrida y, en ese caso, describir de qué tipo
% (por ejemplo: burocracia maquinal-divisional, burocracia profesional simple, etc.).

\endinput
% -------------------------------------------------------------------
% FIN DEL CAPÍTULO — José Ángel Carretero Montes
% -------------------------------------------------------------------
